\documentclass[../main.tex]{subfiles}
\begin{document}

% Chương này chỉ trình bày ý tưởng, không gắn với nút giao hay giờ cụ thể. Hiện thực trên SUMO và số liệu thực
% nghiệm nằm ở Chương 5. Nguồn: PLAN.md §0 (11 hành vi), rules/schema, rules/bounds/param_bounds.yaml,
% src/thesis_sumo/runtime/rules/*.py, src/thesis_sumo/llm/*.

Chương~\ref{chapter:Related_works} cho thấy chưa có hướng nào vừa đưa được hành vi không tuân thủ tại nút vào mô phỏng, vừa chọn được cấu trúc luật cho từng kịch bản. Chương này trình bày tổng quan phương pháp của luận văn (Mục~\ref{section:3.1}) và phần đầu của phương pháp là \emph{sinh luật}. Mục~\ref{section:rg_scenario} định nghĩa kịch bản. Mục~\ref{section:rg_behaviours} trình bày danh mục 11 hành vi được mô hình hóa và cơ sở của từng hành vi. Mục~\ref{section:rg_mechanisms} trình bày cơ sở khoa học và biến điều khiển của từng loại luật. Mục~\ref{section:3.3} định nghĩa đặc tả luật. Mục~\ref{section:rg_llm} trình bày cơ sở và cách thức dùng LLM để sinh luật, và Mục~\ref{section:ev_feedback} mô tả vòng tinh chỉnh theo phản hồi. Chương này chỉ trình bày ý tưởng; phương pháp không gắn với một nút giao hay một khung giờ cụ thể.

\section{Tổng quan phương pháp}
\label{section:3.1}

Mục tiêu của phương pháp là xây dựng một bộ luật hành vi sao cho mô phỏng nút giao với bộ luật đó khớp với dữ liệu quan sát tại nút. Phương pháp gồm hai phần nối thành một vòng lặp (Hình~\ref{fig:1.1}).

Phần \emph{sinh luật} (chương này) xác định thành phần của mô hình hành vi. Đầu vào là mô tả nút giao, danh mục các hành vi và các loại luật, danh mục tham số được phép chỉnh kèm miền giá trị, và phản hồi định lượng của vòng trước. Đầu ra là một \emph{đặc tả luật}: một văn bản có cấu trúc liệt kê các nguyên mẫu tài xế với tỷ trọng và tham số tĩnh, và các luật có điều kiện áp dụng cho từng nguyên mẫu. Đặc tả do LLM sinh, hoặc do người viết tay, hoặc do một bộ tối ưu số điền giá trị; phần thực thi không phân biệt nguồn.

Phần \emph{thực thi luật} (Chương~\ref{chapter:RuleExecution}) gồm mọi bước do mã tất định thực hiện: kiểm tra đặc tả hợp lệ, biên dịch đặc tả thành tham số tĩnh và luật có điều kiện, chạy mô phỏng với một bộ điều khiển đánh giá luật ở mỗi bước, sau đó chuyển kết quả về cùng lược đồ với dữ liệu quan sát và tính các chỉ số. Không có lời gọi LLM nào trong phần này.

Các chỉ số đo được trở thành phản hồi cho lần sinh luật kế tiếp. Vòng lặp dừng khi hết ngân sách số lần chạy mô phỏng, và đặc tả tốt nhất được giữ lại. Cách chia này tạo ra hai giai đoạn thời gian. Giai đoạn \emph{ngoại tuyến} là sinh đặc tả; đây là giai đoạn duy nhất có sự tham gia của LLM. Giai đoạn \emph{trực tuyến} là một lần chạy mô phỏng, hoàn toàn do mã tất định điều khiển, nên lặp lại được với cùng seed.

\section{Định nghĩa kịch bản}
\label{section:rg_scenario}

Một \emph{kịch bản} gồm ba thành phần: một nút giao có đèn, một khung giờ, và dữ liệu quan sát của nút trong khung giờ đó qua nhiều ngày.

Nút giao được mô tả bằng các thông tin đo được tại hiện trường: các nhánh vào và ra, số làn và hướng đi được phép của từng làn, làn rẽ phải liên tục nếu có, chương trình đèn (chu kỳ, pha, thời gian xanh, vàng, đỏ toàn phần), các hướng cấm rẽ, và các không gian mà xe vi phạm có thể dùng (vỉa hè đi được, điểm có thể dừng chờ trong nút). Mô tả này xác định những luật \emph{có thể} áp dụng tại nút: nút không có hướng cấm rẽ trái thì không có rẽ trái trái phép, nút không có vỉa hè đi được thì không có xe đi vỉa hè.

Khung giờ là giờ cao điểm hoặc giờ thấp điểm. Hai khung giờ dùng chung danh mục hành vi và các loại luật, nhưng khác nhau ở ba điểm. Thứ nhất, nhu cầu (số xe theo từng hướng và loại xe trong từng ô 15 phút) lấy từ số đếm của chính khung giờ đó. Thứ hai, trạng thái dòng xe khác nhau: giờ cao điểm có hàng chờ dài và thời gian chờ lâu, giờ thấp điểm có hàng chờ ngắn và khoảng trống lớn hơn giữa các xe. Nhiều luật có điều kiện kích hoạt theo độ dài hàng chờ hoặc thời gian chờ (Mục~\ref{section:rg_mechanisms}), nên cùng một luật tự sinh ra nhiều lượt tìm khoảng trống hơn ở giờ cao điểm mà không cần thay đổi tham số. Thứ ba, tham số hành vi được hiệu chỉnh riêng cho từng khung giờ, vì tỷ trọng người lái vội hay người vi phạm có thể khác nhau giữa hai khung giờ.

Cách đặt kịch bản này cho phép kiểm tra một câu hỏi về hành vi: nếu một đặc tả hiệu chỉnh cho giờ cao điểm cũng tái hiện được giờ thấp điểm chỉ bằng cách thay đổi nhu cầu, thì khác biệt giữa hai khung giờ được giải thích bởi trạng thái dòng xe qua điều kiện kích hoạt; nếu không, tham số hành vi của người tham gia giao thông khác nhau giữa hai khung giờ.

Dữ liệu quan sát của mỗi kịch bản gồm nhiều ngày, chia thành ngày hiệu chỉnh và ngày giữ riêng. Vòng sinh luật chỉ thấy ngày hiệu chỉnh; ngày giữ riêng chỉ dùng để chấm đặc tả cuối cùng, tương tự bước kiểm định trong hiệu chỉnh mô phỏng.

\section{Mô hình hóa hành vi}
\label{section:rg_behaviours}

Danh mục hành vi của luận văn gồm 11 hành vi, xếp theo ba giai đoạn quyết định của một xe khi đi qua nút có đèn. Sáu hành vi có cơ sở trong các nghiên cứu đã công bố; năm hành vi còn lại có cơ sở từ quan sát giao thông thực tế, đối với các hành vi tài liệu chưa đo.

\subsection{Ba giai đoạn quyết định khi đi qua nút}

Mỗi xe đi qua một nút có đèn trải qua ba giai đoạn. Ở giai đoạn \emph{tiếp cận và xếp hàng}, xe đi trên nhánh vào và, nếu đèn đỏ, đứng vào hàng chờ; người lái quyết định vị trí đứng chờ. Ở giai đoạn \emph{tại vạch dừng}, xe đối mặt với tín hiệu đèn; người lái quyết định đi hay chờ. Ở giai đoạn \emph{trong nút}, xe cắt qua các dòng xe khác; người lái quyết định chấp nhận khoảng trống nào và dừng ở đâu. Mô hình chuẩn của trình mô phỏng đã có một hành vi tuân thủ cho mỗi giai đoạn: đi theo làn và xếp hàng sau xe trước, dừng khi đèn đỏ, nhường xe có quyền ưu tiên.

Một hành vi cần mô hình hóa là một cách người tham gia giao thông lệch khỏi hành vi tuân thủ ở một trong ba giai đoạn. Ba giai đoạn cho ba nhóm hành vi: \emph{tìm khoảng trống khi ùn} (giai đoạn tiếp cận và xếp hàng), \emph{không tuân thủ tín hiệu} (giai đoạn tại vạch dừng) và \emph{tương tác trong nút}. Cách phân chia này cung cấp một tiêu chí để rà soát danh mục: với mỗi giai đoạn, danh mục phải chứa các hành vi lệch đã được ghi nhận trong tài liệu về nút có đèn với dòng xe máy chiếm đa số.

\subsection{Hành vi có cơ sở từ tài liệu}

Các nghiên cứu đã công bố cho số đo của sáu hành vi, xếp theo ba nhóm ở trên.

\subsubsection{Không tuân thủ tín hiệu}

Bốn nghiên cứu mô tả bốn khía cạnh của cùng một hành vi. Mai và cộng sự đo trên video tại hai nút có đèn ở Hà Nội: 9,64--11,06\% xe máy vượt đèn đỏ, so với 0,50--0,96\% ô tô~\cite{mai2019signal}. Hai số này cho thấy vượt đèn đỏ là hành vi thường gặp, khác biệt rõ giữa hai loại xe, và chỉ một phần người lái thực hiện. Saval và cộng sự, khi dựng mô hình tác tử cho Hà Nội, cho xe máy ở hàng đầu dừng khi đèn đỏ với xác suất 0,5~\cite{saval2023mixed}; nghĩa là vượt đèn được mô hình hóa như một quyết định xác suất của từng người. Phuengprachuab và cộng sự ghi nhận 61\% số người đi xe máy vượt đỏ tại tám nút ở Thái Lan là vi phạm ``cơ hội'', xuất phát khi thấy ô tô vượt đèn~\cite{phuengprachuab2023opportunistic}. Gao và cộng sự cho thấy nguy cơ vượt đỏ của người đi xe đạp điện ở Thượng Hải tăng 33\% khi thấy người khác vi phạm~\cite{gao2020redlight}. Hai nghiên cứu sau cho thấy quyết định của một người phụ thuộc vào việc người xung quanh vừa làm gì. Điểm chung của bốn nghiên cứu là: vượt đèn đỏ là một quyết định xác suất của từng người (hành vi 1), xác suất này tăng khi có người khác vừa vượt (hành vi 2), và mức vi phạm khác nhau giữa các nhóm người lái. Về xuất phát sớm, Mai và cộng sự đo 13,55--15,82\% xe máy qua vạch từ khoảng 5~s trở lên trước khi đèn xanh~\cite{mai2019signal}, và Phuengprachuab và cộng sự cho thấy việc dừng quá vạch làm tăng xuất phát sớm~\cite{phuengprachuab2023opportunistic}. Từ đó có hành vi 3: xe ở hàng đầu xuất phát trong vài giây cuối của thời gian đỏ.

\subsubsection{Tìm khoảng trống khi ùn}

Chu Cong Minh và cộng sự đo khoảng cách ngang giữa hai xe máy chạy song song ở Hà Nội là 1,10~m~\cite{chu2005speed}. Mohd Zulkiffli và cộng sự ghi nhận hơn 63\% trong 4133 xe máy dừng quá vạch dừng tại bốn nút có đèn ở Malaysia, đo cả giờ cao điểm và giờ thấp điểm~\cite{zulkiffli2023stopping}. Tổng quan của Torquato và cộng sự dẫn các phép đo ở Hà Nội cho thấy thời gian tổn thất khởi động của ô tô tăng từ 3,3~s lên 5,7~s khi có hai hàng xe máy đứng phía trước~\cite{torquato2020motorcycles}, tức xe máy thường lên đứng trước ô tô. Mô hình làn con của Semrau và Erdmann được phát triển để biểu diễn đúng kiểu chiếm không gian này~\cite{semrau2016chinese}. Điểm chung là: khi xếp hàng, xe máy dùng mọi khoảng trống ngang và dọc để lên gần vạch dừng nhất (hành vi 4).

\subsubsection{Tương tác trong nút}

Hsu và cộng sự đo PET trung bình 2,15~s ở xung đột rẽ trái cắt ngang có xe máy tại Đài Loan~\cite{hsu2024v2x}; PET nhỏ nghĩa là người tham gia giao thông chấp nhận khoảng trống nhỏ khi cắt qua dòng xe khác (hành vi 9). Mai và cộng sự mô tả xe máy của luồng phụ đi vào giữa nút rồi dừng chờ khi pha đèn chuyển~\cite{mai2019signal}; đây là cơ sở định tính cho hành vi 10, dùng khoảng trống trong nút làm vị trí chờ để rẽ.

Các nghiên cứu này đo ở địa điểm, thời điểm và với định nghĩa khác nhau, và một số không ở Việt Nam. Các nghiên cứu này đủ để xác định \emph{sự tồn tại} và \emph{cơ chế} của hành vi, nhưng chưa đủ để xác định giá trị tham số của một nút cụ thể. Giá trị tham số vì vậy được hiệu chỉnh từ dữ liệu quan sát của từng kịch bản.

\subsection{Hành vi có cơ sở từ quan sát}

Các tài liệu đã khảo sát mô tả xe máy dùng khoảng trống \emph{trong làn} của mình (hành vi 4), nhưng chưa đo trường hợp hàng chờ dài tới mức xe tìm không gian \emph{ngoài làn}: đi sai làn để lên đầu hàng rồi cắt về (hành vi 5), dừng chờ đèn đỏ trên làn rẽ phải liên tục và chặn xe rẽ phải (hành vi 6), đi trên vỉa hè để vượt hàng chờ (hành vi 7), và đi ngược chiều ở đoạn không có dải phân cách (hành vi 8). Luận văn bổ sung bốn hành vi này từ quan sát giao thông thực tế của tác giả, như các phần mở rộng của cùng cơ chế ở hành vi 4 sang không gian khác. Hành vi xe máy đi chậm giữa làn (hành vi 11) cũng đến từ quan sát. Quan sát cũng cho thấy phần lớn xe máy đứng sau một xe đang chờ quay đầu sẽ lách qua xe đó, còn ô tô ít khi làm vậy; luận văn thêm một luật bổ sung cho tình huống này. Luận văn chưa tìm thấy nghiên cứu nào đo các hành vi bổ sung này, nên giá trị mục tiêu của chúng mang nhãn giả định cho tới khi có dữ liệu video.

\begin{table}[!htbp]
\centering
\small
\caption{Mười một hành vi trong danh mục và cơ sở của từng hành vi.}
\label{table:1.1}
\begin{tabular}{|c|p{2.6cm}|p{5.2cm}|p{4.6cm}|}
\hline
\textbf{STT} & \textbf{Nhóm} & \textbf{Hành vi} & \textbf{Cơ sở} \\
\hline
1 & \multirow{3}{2.6cm}{Không tuân thủ tín hiệu} & Ngưỡng kiên nhẫn và vượt đèn đỏ & \cite{mai2019signal,saval2023mixed} \\
2 & & Vượt đèn đỏ theo đoàn & \cite{phuengprachuab2023opportunistic,gao2020redlight} \\
3 & & Xuất phát sớm trước khi đèn xanh & \cite{mai2019signal,phuengprachuab2023opportunistic} \\
\hline
4 & \multirow{5}{2.6cm}{Tìm khoảng trống khi ùn} & Khoảng cách ngang nhỏ, lấn chiếm, luồn lách, dừng quá vạch & \cite{chu2005speed,zulkiffli2023stopping,torquato2020motorcycles} \\
5 & & Đi sai làn để lên đầu hàng rồi cắt về & Quan sát \\
6 & & Dừng chờ đèn đỏ trên làn rẽ phải liên tục & Quan sát \\
7 & & Xe máy đi trên vỉa hè để vượt hàng chờ & Quan sát \\
8 & & Đi ngược chiều ở đoạn không có dải phân cách & Quan sát \\
\hline
9 & \multirow{3}{2.6cm}{Tương tác trong nút} & Chấp nhận khoảng trống nhỏ, tạt đầu & \cite{hsu2024v2x} (gián tiếp) \\
10 & & Rẽ trái trái phép: vào giữa nút, dừng chờ, rẽ khi pha đi thẳng đối diện kết thúc & \cite{mai2019signal} (mô tả định tính) \\
11 & & Xe máy đi chậm giữa làn & Quan sát \\
\hline
\end{tabular}
\end{table}

\subsection{Độ tin cậy của danh mục}

Danh mục 11 hành vi (Bảng~\ref{table:1.1}) dựa trên ba lập luận. Thứ nhất, mỗi giai đoạn quyết định có ít nhất một hành vi rút từ tài liệu, và mỗi hành vi rút từ tài liệu có một đại lượng đo được trong tài liệu (tỷ lệ vượt đỏ, tỷ lệ xuất phát sớm, khoảng cách ngang, tỷ lệ dừng quá vạch, PET), nên có thể đặt giá trị mục tiêu cho nó. Riêng PET của hành vi 9 chưa có chỉ số tương ứng trong luận văn (Mục~\ref{section:3.7}). Thứ hai, mỗi hành vi bổ sung từ quan sát được gắn rõ nhãn và không được dùng như số liệu có nguồn; khi tính điểm tổng, phần sai lệch từ các mục tiêu giả định được báo cáo tách riêng (Mục~\ref{section:3.7}). Thứ ba, luận văn không khẳng định danh mục là đầy đủ. Hai chỉ số đo khả năng thông hành và thời gian đi qua nút chịu ảnh hưởng của mọi hành vi; nếu danh mục thiếu một hành vi quan trọng, sai lệch sẽ còn lại ở hai chỉ số này dù mọi chỉ số riêng của từng hành vi đã khớp. Danh mục vì vậy kiểm chứng được bằng dữ liệu.

Có hai giới hạn phạm vi. Trong một nút có làn rẽ phải liên tục, rẽ phải khi đèn đỏ là đúng luật; phương pháp không coi đó là vi phạm. Một hành vi cần không gian riêng chỉ được đưa vào mô hình nếu qua thí nghiệm khả thi mô tả ở Mục~\ref{section:ex_space}; hành vi không qua thí nghiệm vẫn thuộc danh mục nhưng nằm ngoài phạm vi mô hình.

\section{Cơ sở khoa học của các loại luật và biến điều khiển}
\label{section:rg_mechanisms}

Mỗi hành vi được biểu diễn theo một trong hai cách. Hành vi xảy ra liên tục và không phụ thuộc tình huống, như giữ khoảng cách ngang nhỏ, được biểu diễn bằng \emph{tham số tĩnh} của nguyên mẫu tài xế. Hành vi chỉ xảy ra trong một tình huống, như vượt đèn đỏ khi vừa có người vượt, được biểu diễn bằng \emph{luật có điều kiện}, gồm điều kiện kích hoạt, hành động và các tham số. Mỗi loại luật dưới đây được trình bày theo ba phần: cơ sở của nó trong tài liệu, dạng toán học hoặc logic, và các biến điều khiển.

Ngoài tham số mô tả bản thân hành vi, nhiều luật có một \emph{biến điều khiển trung gian} (gọi tắt là biến đệm): một điều kiện hoặc giới hạn đặt giữa điều kiện kích hoạt và hành động. Biến đệm không tạo ra hành vi; nó ngăn hành vi gây ra hậu quả mà thực tế không có. Mỗi biến đệm được nêu cùng hậu quả nếu thiếu nó.

\subsection{Nguyên mẫu tài xế và tham số tĩnh}

Số liệu của Mai và cộng sự cho thấy chỉ khoảng một phần mười số xe máy vượt đèn đỏ, và tỷ lệ này ở ô tô thấp hơn khoảng mười lần~\cite{mai2019signal}; Saval và cộng sự hiệu chỉnh tham số riêng cho ô tô và xe máy~\cite{saval2023mixed}. Hai kết quả này cho thấy người tham gia giao thông không đồng nhất: trong cùng một dòng xe máy có người tuân thủ và có người vi phạm. Phương pháp biểu diễn điều đó bằng các \emph{nguyên mẫu tài xế}. Mỗi loại phương tiện được chia thành vài nguyên mẫu, mỗi nguyên mẫu có tỷ trọng trong dòng xe và tập tham số tĩnh riêng thuộc bốn nhóm: bám xe (giãn cách thời gian mong muốn tới xe trước, khoảng cách tối thiểu khi dừng, gia tốc, giảm tốc, mức dao động tốc độ), di chuyển ngang (khoảng cách ngang tối thiểu, tốc độ ngang), chuyển làn (mức chen lấn, mức mất kiên nhẫn) và ứng xử tại nút (khoảng trống chấp nhận, mức sẵn sàng đi khi đèn vừa đỏ). Hành vi 4 và 9 nằm chủ yếu ở các tham số này: khoảng cách ngang tối thiểu tương ứng với số đo 1,10~m của Chu Cong Minh và cộng sự~\cite{chu2005speed}, khoảng trống chấp nhận tương ứng với PET nhỏ của Hsu và cộng sự~\cite{hsu2024v2x}.

Xe vi phạm thuộc một \emph{lớp vi phạm} riêng. Biến đệm ở đây là quyền dùng không gian: chỉ lớp vi phạm được vào những không gian xe tuân thủ không được vào, như vỉa hè hay hướng rẽ bị cấm. Nếu không tách lớp, mở một không gian cho hành vi vi phạm đồng nghĩa với mở nó cho mọi xe, và nguyên mẫu tuân thủ cũng sẽ dùng vỉa hè.

\subsection{Vượt đèn đỏ có hiệu ứng bầy đàn}

Hành vi 1 và 2 dùng chung một luật, xây dựng từ điểm chung của bốn nghiên cứu ở Mục~\ref{section:rg_behaviours}: quyết định xác suất của từng người~\cite{saval2023mixed,mai2019signal} và ảnh hưởng của người xung quanh~\cite{phuengprachuab2023opportunistic,gao2020redlight}. Khi xe đến vạch dừng lúc đèn đỏ, xác suất xe quyết định vượt là
\begin{equation}
p = \min\left(p_0 + k\,n,\; p_{\max}\right),
\label{eq:swarm}
\end{equation}
trong đó $p_0$ là xác suất vượt đỏ nền (ngưỡng kiên nhẫn, hành vi 1), $n$ là số xe đã vượt đỏ trên cùng nhánh trong $\Delta t$ giây gần nhất, và $k$ là mức tăng xác suất trên mỗi xe vừa vượt (hiệu ứng bầy đàn, hành vi 2). Dạng cộng tuyến tính là dạng đơn giản nhất có hai tính chất mà tài liệu đòi hỏi: khi không ai vừa vượt, xác suất bằng mức nền; mỗi người vừa vượt làm xác suất tăng thêm. Nếu thiếu $n$, các lượt vượt đỏ độc lập với nhau, và mô hình không tái hiện được việc vượt đỏ dồn theo đợt (phương sai số xe vượt đỏ mỗi chu kỳ lớn hơn trung bình).

Luật có hai biến đệm. Biến thứ nhất là trần $p_{\max}$: nếu thiếu trần, mỗi lượt vượt làm tăng xác suất của lượt sau, và một chuỗi vượt đỏ có thể kéo tới $p = 1$ cho cả nhánh. Biến thứ hai là \emph{điều kiện nút thông}: xe đã quyết định vượt chỉ vào nút khi làn ra còn chỗ, không có xe trên các luồng cắt ngang và không có xe đứng trên chính luồng của nó trong nút. Nếu nút chưa thông, xe đợi ở vạch, và nếu đèn hết đỏ trước khi nút thông thì lượt đó không xảy ra. Nếu thiếu điều kiện này, xe vượt đỏ có thể dừng lại giữa dòng xe đang có đèn xanh, chặn nút và gây khóa nút, trong khi người vượt đèn trên thực tế tránh hậu quả này vì chỉ đi khi thấy đường thông.

\subsection{Xuất phát sớm}

Hành vi 3 dựa trên số đo xuất phát sớm của Mai và cộng sự~\cite{mai2019signal}. Luật được kích hoạt trong $t_{\text{before}}$ giây cuối của thời gian đỏ: xe đứng đầu hàng nhích qua vạch dừng với xác suất $p$. Biến đệm là tốc độ tối đa $v_{\max}$ khi nhích. Nếu thiếu giới hạn này, xe xuất phát sớm tăng tốc như khi có đèn xanh và đi vào luồng cắt ngang chưa xả hết; với giới hạn này, xe chỉ nhích dần lên phía trước như người xuất phát sớm trên thực tế, và các cơ chế nhường đường vẫn hoạt động.

\subsection{Tìm khoảng trống khi ùn}

Hành vi 5, 6 và 7 dùng chung một luật vì chung một nguyên nhân: hàng chờ dài. Do chưa có tài liệu đo các hành vi này, điều kiện kích hoạt theo độ dài hàng chờ và thời gian chờ là giả thuyết của luận văn, cùng dạng với ngưỡng kiên nhẫn ở hành vi 1: người lái chấp nhận chờ tới một ngưỡng, vượt ngưỡng đó thì tìm phương án khác. Mỗi xe thuộc nguyên mẫu tìm khoảng trống được \emph{rút ngẫu nhiên đúng một lần}, tại thời điểm đầu tiên một trong ba điều kiện kích hoạt thỏa: hàng chờ ở làn đúng của xe dài ít nhất $q$ mét; xe đã dừng ít nhất $t_w$ giây; hoặc, $t_g$ giây trước lượt xanh kế tiếp, xe ước lượng rằng mình không qua kịp vạch trong lượt xanh đó. Kết quả rút là một trong các lựa chọn: đi làn sai rồi cắt về làn đúng gần vạch dừng, dừng trên làn rẽ phải liên tục, lên vỉa hè, hoặc không làm gì. Xe đi trong không gian thay thế với tốc độ không quá $v_{\text{creep}}$.

Luật có ba biến đệm. Thứ nhất là \emph{rút một lần}: nếu xe được rút lại ở mỗi bước mô phỏng, xác suất làm hành vi trong một lần chờ dài tiến tới 1 bất kể xác suất mỗi bước nhỏ đến đâu, và tham số xác suất mất ý nghĩa. Thứ hai là \emph{sai số ước lượng} $\varepsilon$: xe ước lượng thời gian cần để qua vạch với một sai số nhân ngẫu nhiên. Nếu ước lượng chính xác, mọi xe ở cùng vị trí trong hàng quyết định giống nhau, trong khi người thật đánh giá hàng chờ khác nhau. Thứ ba là \emph{vùng nhập lại} và $v_{\text{creep}}$: xe đi sai làn phải cắt về làn đúng trong một đoạn gần vạch dừng, và nếu chưa nhập được thì đứng ở cuối làn sai, chặn xe sau. Nếu thiếu điều này, xe đi sai làn sẽ đi thẳng qua nút từ làn không có hướng của nó.

Xác suất của một lựa chọn không có không gian tại nút đang xét (ví dụ nút không có vỉa hè đi được) được chia lại theo tỷ lệ cho các lựa chọn còn lại, để cùng một bộ luật dùng được ở nút khác. Vì điều kiện kích hoạt phụ thuộc trạng thái hàng chờ, luật này cho số lượt khác nhau giữa giờ cao điểm và giờ thấp điểm với cùng tham số (Mục~\ref{section:rg_scenario}).

\subsection{Rẽ trái trái phép chờ giữa nút}

Hành vi 10 dựa trên mô tả định tính của Mai và cộng sự về xe máy luồng phụ vào giữa nút rồi dừng chờ~\cite{mai2019signal}, và gồm ba bước: xe đi vào giữa nút trong pha xanh đi thẳng của nhánh mình, dừng chờ ở một điểm trong nút, rồi rẽ khi dòng đi thẳng đối diện ngớt. Không gian cho hướng rẽ bị cấm là một bản vá mạng chỉ mở cho lớp vi phạm. Việc xe rời điểm chờ qua khoảng trống của dòng đối diện dùng cơ chế nhường đường sẵn có của trình mô phỏng, điều khiển bằng tham số khoảng trống chấp nhận của lớp vi phạm; đây là cách mô hình hóa chấp nhận khoảng trống quen thuộc trong mô phỏng vi mô. Luật có điều kiện chỉ thêm một tham số: độ trễ giải phóng $\Delta_{\text{release}}$, khoảng thời gian giữ xe thêm sau lúc trình mô phỏng cho phép xe đi.

Biến đệm của hành vi này là nguyên tắc \emph{chỉ giữ, không đẩy}: luật được phép giữ xe lâu hơn, không bao giờ được cho xe đi xuyên qua dòng xung đột. Một cách làm trực tiếp là cho các xe cắt ngang bỏ qua xe đang chờ, hoặc buộc xe chờ đi đúng thời điểm mong muốn. Cách này làm trình mô phỏng không còn tính xung đột giữa các xe đó, nên chúng đi xuyên qua nhau. Cái giá của nguyên tắc là luật chỉ kéo dài được thời gian giải phóng vùng chờ, không rút ngắn được.

\subsection{Đi chậm giữa làn và vượt xe quay đầu}

Hành vi 11 được biểu diễn bằng một luật giữ xe máy đi chậm ở một vị trí ngang lệch $\delta$ so với tim làn. Luật bổ sung cho xe chờ quay đầu được kích hoạt khi xe đứng sau một xe đang chờ quay đầu trên cùng làn lúc đèn của chính nó xanh. Sau $t_{\text{wait}}$ giây, mỗi cặp xe được rút một lần với xác suất $p$; nếu trúng, xe lách ngang vào khe gần nhất bên cạnh xe quay đầu, hoặc đổi sang làn kề có hướng đi của nó. Biến đệm là thời gian chờ $t_{\text{wait}}$ và quy tắc rút một lần cho mỗi cặp: nếu thiếu, xe lách ngay khi xe phía trước vừa dừng, kể cả khi xe đó chỉ dừng trong chốc lát.

\subsection{Tổng hợp các loại luật}

Bảng~\ref{table:rg_mech} tóm tắt các loại luật, hành vi tương ứng, tham số và biến đệm. Miền giá trị cụ thể của từng tham số phụ thuộc trình mô phỏng và nằm ở Chương~\ref{chapter:Experiment} (Bảng~\ref{table:3.2}).

\begin{table}[!htbp]
\centering
\small
\caption{Các loại luật có điều kiện, tham số và biến đệm.}
\label{table:rg_mech}
\begin{tabular}{|p{2.7cm}|p{1.3cm}|p{4.0cm}|p{5.2cm}|}
\hline
\textbf{Luật} & \textbf{Hành vi} & \textbf{Tham số hành vi} & \textbf{Biến đệm và hậu quả nếu thiếu} \\ \hline
\texttt{swarm\_\allowbreak red\_\allowbreak running} & 1, 2 & $p_0$, $k$, $\Delta t$ & Trần $p_{\max}$ (thiếu: chuỗi vượt đỏ tự khuếch đại); điều kiện nút thông (thiếu: khóa nút) \\ \hline
\texttt{early\_\allowbreak start} & 3 & $t_{\text{before}}$, $p$ & Tốc độ nhích $v_{\max}$ (thiếu: xe đi thẳng vào luồng cắt ngang) \\ \hline
\texttt{queue\_\allowbreak space\_\allowbreak seeking} & 5, 6, 7 & $q$, $t_w$, $t_g$, ba xác suất lựa chọn & Rút một lần (thiếu: xác suất tiến tới 1 khi chờ lâu); sai số $\varepsilon$ (thiếu: mọi xe quyết định như nhau); vùng nhập lại, $v_{\text{creep}}$ (thiếu: đi qua nút từ làn sai) \\ \hline
\texttt{illegal\_\allowbreak left\_\allowbreak box} & 10 & $\Delta_{\text{release}}$ & Chỉ giữ, không đẩy (thiếu: xe đi xuyên qua nhau) \\ \hline
\texttt{slow\_\allowbreak lane\_\allowbreak center} & 11 & $\delta$ & -- \\ \hline
\texttt{uturn\_\allowbreak bypass} & bổ sung & $p$ & Thời gian chờ $t_{\text{wait}}$, rút một lần mỗi cặp (thiếu: lách ngay cả khi xe trước chỉ dừng chốc lát) \\ \hline
\end{tabular}
\end{table}

\section{Đặc tả luật}
\label{section:3.3}

Một đặc tả luật gồm phần định danh, phần truy vết và hai thành phần chính (Bảng~\ref{table:3.1a}). Đặc tả là văn bản có cấu trúc theo một lược đồ cố định, nên kiểm tra được bằng máy và so sánh được giữa các phương án.

\begin{table}[!htbp]
\centering
\small
\caption{Cấu trúc của một đặc tả luật.}
\label{table:3.1a}
\begin{tabular}{|p{3.4cm}|p{10.6cm}|}
\hline
\textbf{Thành phần} & \textbf{Nội dung} \\
\hline
Định danh, nguồn gốc & Mã đặc tả; nguồn gốc là thủ công, LLM hoặc bộ tối ưu số \\
\hline
Truy vết & Mô hình đã sinh đặc tả, phiên bản prompt, vòng lặp, chỉ số mẫu \\
\hline
Nguyên mẫu tài xế & Danh sách nguyên mẫu. Mỗi nguyên mẫu có: loại phương tiện cơ sở (ô tô, xe máy, hoặc lớp vi phạm của chúng); phân phối mà nó thuộc về; tỷ trọng trong phân phối đó; tập tham số tĩnh; lý do chọn (bắt buộc với đặc tả do LLM sinh). Tổng tỷ trọng trong mỗi phân phối bằng 1 \\
\hline
Luật có điều kiện & Danh sách luật. Mỗi luật có loại luật, các nguyên mẫu áp dụng và tham số \\
\hline
\end{tabular}
\end{table}

Không gian tham số nằm trong một \emph{tệp miền giá trị} dùng chung cho mọi nguồn sinh đặc tả. Mỗi dòng ghi tên tham số, đơn vị, cận dưới, cận trên, giá trị mặc định, nhóm hành vi liên quan, các chỉ số có thể ràng buộc tham số, và nguồn của khoảng giá trị kèm nhãn bằng chứng (tài liệu của trình mô phỏng, tài liệu khoa học, chuyển giao từ bối cảnh khác, quyết định thiết kế, hoặc giả định). Hai quy tắc giữ không gian này nhỏ và có nghĩa. Tham số không chỉ số nào ràng buộc được cố định ở giá trị mặc định, vì không có cách hiệu chỉnh nó từ dữ liệu. Tham số cho phép xe bỏ qua xe xung đột có quyền ưu tiên bị loại hoàn toàn, vì theo thiết kế nó tạo ra va chạm thật. Dùng chung một tệp miền giá trị cũng là điều kiện để so sánh LLM với bộ tối ưu số: hai bên tìm trên đúng cùng một tập.

\section{Sinh luật bằng LLM}
\label{section:rg_llm}

\subsection{Cơ sở của việc dùng LLM sinh luật}

Hiệu chỉnh một mô hình mô phỏng vi mô là chỉnh mô hình cho tới khi đầu ra khớp số đo, lặp lại qua nhiều lần chạy~\cite{hollander2008calibration}. Trong bài toán của luận văn, việc chỉnh có hai phần: chọn \emph{cấu trúc} (cần những nguyên mẫu nào, luật nào áp dụng cho nguyên mẫu nào) và chọn \emph{giá trị} tham số. Bộ tối ưu số chỉ thực hiện được phần thứ hai, và cần cấu trúc có sẵn~\cite{akiba2019optuna}. Phần thứ nhất là một lựa chọn rời rạc trên tổ hợp các loại luật và nguyên mẫu, phụ thuộc vào mô tả nút giao, và hiện do người lập mô hình làm bằng hiểu biết về hành vi người lái.

Ba nhóm kết quả cho thấy LLM có thể đảm nhận phần này. Thứ nhất, LLM có thể hoạt động như một bộ tối ưu khi bài toán được mô tả bằng ngôn ngữ tự nhiên: trong OPRO, mỗi bước LLM nhận các lời giải trước cùng điểm của chúng và đề xuất lời giải mới, và prompt mà OPRO tìm được cho kết quả cao hơn prompt do người viết, tới 8\% trên GSM8K và tới 50\% trên Big-Bench Hard~\cite{yang2024opro}. Vòng sinh luật của luận văn có đúng dạng đó: prompt chứa các đặc tả trước cùng sai lệch của chúng. Thứ hai, LLM sinh được đầu ra có cấu trúc mà một bộ đánh giá tự động chấm được, và việc ghép LLM với bộ đánh giá trong một vòng lặp đã tìm được lời giải mới: FunSearch ghép LLM với bộ đánh giá chương trình và tìm được các cấu trúc cap set lớn hơn cấu trúc tốt nhất đã biết~\cite{romeraparedes2024funsearch}; Eureka để LLM viết mã hàm thưởng cho học tăng cường, chấm bằng mô phỏng, và sửa theo kết quả mô phỏng qua nhiều vòng~\cite{ma2024eureka}. Điểm chung của hai công trình là đề xuất của LLM không được chấp nhận trực tiếp: mọi đề xuất phải qua bộ đánh giá, nên lỗi của LLM bị loại trước khi đi vào kết quả. Thứ ba, LLM đọc được mô tả nút giao bằng ngôn ngữ tự nhiên cùng bảng sai lệch định lượng, và trả lời kèm lý do; lý do này cho người đọc kiểm tra một lựa chọn cấu trúc có hợp với hiểu biết về hành vi hay không, điều bộ tối ưu số không cung cấp. Các công trình gần đây đã dùng LLM để sinh và chỉnh kịch bản mô phỏng giao thông từ yêu cầu bằng ngôn ngữ tự nhiên~\cite{jeong2025agentsumo,li2026chatsumo}; luận văn dùng cùng khả năng đó cho luật hành vi.

Dùng LLM cũng có rủi ro: đầu ra có thể sai lược đồ, dùng tên tham số không tồn tại, hoặc chọn giá trị ngoài miền. Theo đúng điểm chung của FunSearch và Eureka, tính đúng đắn của phương pháp không dựa vào LLM mà dựa vào bộ lọc phía sau, với ba ràng buộc. LLM chỉ chọn trong danh mục tham số và loại luật có sẵn, với miền giá trị cho trước. Mọi đặc tả qua bộ kiểm tra hợp lệ trước khi chạy (Mục~\ref{section:3.4}), và mọi đặc tả hợp lệ được chấm bằng mô phỏng trên dữ liệu quan sát. LLM không bao giờ điều khiển xe trực tiếp.

\subsection{Đầu vào và đầu ra của LLM}

Prompt gồm ba phần: vai trò và các ràng buộc; bối cảnh nút giao kèm từ điển tham số sinh tự động từ tệp miền giá trị và giá trị mục tiêu của các ngày hiệu chỉnh; và với mỗi vòng trước, các đặc tả của vòng đó cùng phản hồi định lượng (Mục~\ref{section:ev_feedback}). Đầu ra là một đặc tả theo lược đồ của Bảng~\ref{table:3.1a}. Prompt có phiên bản và được khóa theo từng chiến dịch, để mọi kết quả gắn với đúng prompt đã tạo ra nó.

\subsection{Hai chế độ sinh luật}

Để tách hai vai trò của LLM, vòng lặp có hai chế độ. Chế độ \emph{chỉ tham số} khóa cấu trúc của một bộ luật viết tay; LLM chỉ thay đổi giá trị, và mọi thay đổi cấu trúc bị từ chối. Ở chế độ này LLM làm cùng việc với bộ tối ưu số, nên hai bên so sánh trực tiếp được. Chế độ \emph{đầy đủ} cho LLM sinh cả cấu trúc lẫn tham số. Chế độ này có một rủi ro: LLM có thể bỏ bớt luật, tức bỏ bớt hành vi mà mô hình phải tái hiện. Một hành vi bị bỏ không có sự kiện nào trong mô phỏng, và với tỷ lệ nhỏ, khoảng cách của nó có thể nhỏ hơn khoảng cách của một hành vi xảy ra quá mức (Mục~\ref{section:3.7}).

\section{Tinh chỉnh lặp theo phản hồi}
\label{section:ev_feedback}

Mỗi vòng sinh một số mẫu đặc tả độc lập, mỗi mẫu dùng một seed riêng để các mẫu khác nhau. Một mẫu không hợp lệ được gửi lại cho LLM kèm danh sách lỗi để sửa, với số lần thử giới hạn; đặc tả không hợp lệ không tiêu tốn lần chạy mô phỏng nào. Các mẫu hợp lệ được thực thi và chấm theo Chương~\ref{chapter:RuleExecution}.

Sau khi các đặc tả của một vòng được chấm, phản hồi cho LLM được tạo từ bốn nguồn: khoảng cách và tỷ số $R$ theo từng chỉ số (Mục~\ref{section:3.7}), kèm giá trị quan sát và giá trị mô phỏng của từng đặc tả; tỷ lệ ô qua cổng lọc lưu lượng; số và vị trí xe bị kẹt; và tóm tắt nhật ký can thiệp (số lần mỗi luật được kích hoạt và số lần trình mô phỏng thi hành). Lỗi kiểm tra hợp lệ, nếu có, được gửi kèm. Vòng gần nhất được gửi đầy đủ; các vòng cũ hơn chỉ còn vài dòng tóm tắt, để prompt không vượt giới hạn ngữ cảnh của LLM.

Phản hồi có một biến đệm về dữ liệu: nó chỉ dùng các ngày hiệu chỉnh, và mã báo lỗi nếu tên một ngày giữ riêng xuất hiện trong văn bản phản hồi. Nếu thiếu ràng buộc này, LLM thấy được dữ liệu dùng để chấm đặc tả cuối cùng, và điểm trên ngày giữ riêng không còn đo khả năng tổng quát hóa.

Mọi nguồn sinh đặc tả dùng chung một \emph{ngân sách} số lần chạy mô phỏng, gắn với từng chiến dịch và lưu trên đĩa, nên chạy lại hay bị ngắt giữa chừng đều không làm vượt ngân sách. Khi hết ngân sách, đặc tả tốt nhất trong số các đặc tả đạt tiêu chí chấp nhận (Mục~\ref{section:ev_eligibility}) được giữ lại. Mọi yêu cầu gửi LLM và câu trả lời được lưu nguyên văn; khi chạy lại với cùng tin nhắn và cùng định danh mô hình, câu trả lời đã lưu được phát lại, nên một chiến dịch tái lập được mà không cần gọi lại LLM.

\medskip

Chương này đã nêu tổng quan phương pháp và trình bày phần sinh luật: định nghĩa kịch bản, danh mục hành vi có cơ sở từ tài liệu và quan sát, cơ sở và biến đệm của từng loại luật, đặc tả luật, lý do và cách dùng LLM, và vòng tinh chỉnh theo phản hồi. Một đặc tả do LLM sinh mới chỉ là văn bản. Chương~\ref{chapter:RuleExecution} trình bày cách kiểm tra nó, biến nó thành hành vi của xe trong mô phỏng mà không làm hỏng các cơ chế an toàn của trình mô phỏng, và đo kết quả.

\end{document}
